\section{Características de los Usuarios}

Se identifican cuatro perfiles de usuario que interactúan con el framework Kit Pagos Colombia, diferenciados por su nivel de experiencia técnica, su frecuencia de uso y su relación con los componentes del sistema.

\subsection{Desarrollador de software junior}

\begin{table}[H]
  \centering
  \renewcommand{\arraystretch}{1.2}
  \begin{tabularx}{\textwidth}{|l|X|}
    \hline
    \textbf{Característica} & \textbf{Descripción} \\
    \hline
    Descripción & Desarrollador con menos de dos años de experiencia profesional que se enfrenta por primera vez a la integración de una pasarela de pago. Puede ser estudiante universitario, freelancer o parte de un equipo de una startup en etapa temprana. \\
    \hline
    Experiencia técnica & Conocimiento básico de REST APIs y JavaScript/TypeScript. Poca o nula experiencia previa con pasarelas de pago colombianas. \\
    \hline
    Frecuencia de uso & Esporádica. Utiliza el framework durante el desarrollo inicial de un proyecto y ocasionalmente para agregar nuevas pasarelas. \\
    \hline
    Necesidades principales & Documentación clara, ejemplos de código funcionales, datos de prueba accesibles y un entorno de simulación que no requiera configuración compleja. \\
    \hline
    Interacción con el framework & Principalmente con el SDK de servidor, el módulo de navegador para tokenizar la tarjeta, la documentación de datos de prueba y la API de Simulación en local. \\
    \hline
  \end{tabularx}
  \caption{Perfil del desarrollador junior.}
  \label{tab:perfil-junior}
\end{table}

\subsection{Desarrollador de software senior}

\begin{table}[H]
  \centering
  \renewcommand{\arraystretch}{1.2}
  \begin{tabularx}{\textwidth}{|l|X|}
    \hline
    \textbf{Característica} & \textbf{Descripción} \\
    \hline
    Descripción & Desarrollador con experiencia en integraciones de sistemas y arquitecturas de software. Ha trabajado previamente con APIs de pago o servicios externos. Puede ser parte de un equipo de producto o consultora tecnológica. \\
    \hline
    Experiencia técnica & Dominio avanzado de TypeScript, Node.js, patrones de diseño y arquitecturas de integración. Familiaridad con los conceptos de autenticación, webhooks y manejo de errores en APIs externas. \\
    \hline
    Frecuencia de uso & Regular. Utiliza el framework como parte de proyectos activos y puede contribuir a su extensión con nuevos adaptadores. \\
    \hline
    Necesidades principales & Flexibilidad arquitectónica, capacidad de extensión para agregar nuevas pasarelas, control sobre el comportamiento de reintentos y manejo avanzado de errores. \\
    \hline
    Interacción con el framework & Con todos los componentes del framework, incluyendo la implementación de adaptadores personalizados, la configuración avanzada del SDK y la capa REST \texttt{/v1/api}. \\
    \hline
  \end{tabularx}
  \caption{Perfil del desarrollador senior.}
  \label{tab:perfil-senior}
\end{table}

\subsection{Desarrollador frontend}

\begin{table}[H]
  \centering
  \renewcommand{\arraystretch}{1.2}
  \begin{tabularx}{\textwidth}{|l|X|}
    \hline
    \textbf{Característica} & \textbf{Descripción} \\
    \hline
    Descripción & Desarrollador que construye la página de pago del comercio. Trabaja en el navegador y puede no conocer los detalles de autenticación de cada pasarela. \\
    \hline
    Experiencia técnica & Dominio de JavaScript/TypeScript, del DOM y de peticiones HTTP desde el navegador. Familiaridad con formularios de tarjeta y con Content-Security-Policy. \\
    \hline
    Frecuencia de uso & Durante la implementación del checkout y cada vez que se agrega una pasarela o se ajusta la página de pago. \\
    \hline
    Necesidades principales & Un módulo de tokenización que no exponga el número de tarjeta al servidor del comercio, una API REST para cobrar sin instalar el SDK, y ejemplos de protección de la página. \\
    \hline
    Interacción con el framework & Principalmente con \texttt{kit-pagos-colombia/browser} y con \texttt{/v1/api}; consulta la documentación de datos de prueba y la guía de implementación para la Content-Security-Policy. \\
    \hline
  \end{tabularx}
  \caption{Perfil del desarrollador frontend.}
  \label{tab:perfil-frontend}
\end{table}

\subsection{Evaluador / Tester de software}

\begin{table}[H]
  \centering
  \renewcommand{\arraystretch}{1.2}
  \begin{tabularx}{\textwidth}{|l|X|}
    \hline
    \textbf{Característica} & \textbf{Descripción} \\
    \hline
    Descripción & Persona responsable de verificar el correcto funcionamiento de la integración de pagos en una aplicación. Puede ser un QA engineer, un desarrollador en fase de pruebas o un integrante del equipo de proyecto encargado de validar los flujos de pago antes de salir a producción. \\
    \hline
    Experiencia técnica & Conocimiento funcional de los flujos de pago. Experiencia con herramientas de prueba de APIs (Postman, curl). No necesariamente escribe código de integración. \\
    \hline
    Frecuencia de uso & Intensiva durante los ciclos de prueba del proyecto. Usa principalmente la API de Simulación y la documentación de datos de prueba. \\
    \hline
    Necesidades principales & Acceso a escenarios de prueba predefinidos (éxito, rechazo, error de red, timeout), datos de prueba válidos por pasarela y una API de simulación estable y predecible. \\
    \hline
    Interacción con el framework & Principalmente con \texttt{/v1/sim} y \texttt{/v1/api} a través de herramientas REST, con la colección de Postman y con la documentación centralizada de datos de prueba. \\
    \hline
  \end{tabularx}
  \caption{Perfil del evaluador/tester.}
  \label{tab:perfil-tester}
\end{table}
